% year: תשע"ט
% semester: ב
% moed: א
% number: 1ג
% points: 5
% categories: definite-integrals, derivatives
% source: exams/מבחנים/תשעט/חודיסמן מועד א תשעט.pdf

\begin{question}
(5 נק') מהו שיפוע המשיק לקו $y=\int_0^{x+1}e^{-2t^2}\,dt$ בנקודה עליו $x=0$.
\end{question}

\begin{hint}
השתמשו במשפט היסודי של החשבון האינטגרלי יחד עם כלל השרשרת.
\end{hint}

\begin{solution}
נסמן $F(u)=\int_0^u e^{-2t^2}\,dt$. הפונקציה $e^{-2t^2}$ רציפה, ולכן לפי המשפט היסודי $F'(u)=e^{-2u^2}$. מכאן $y(x)=F(x+1)$ ולפי כלל השרשרת
\[ y'(x)=F'(x+1)\cdot 1=e^{-2(x+1)^2}. \]
שיפוע המשיק בנקודה $x=0$ הוא $y'(0)=e^{-2}=\dfrac{1}{e^2}$.
\end{solution}
